\subsection{Plataforma, seguridad y servicios comunes}
Las cuentas 1.2, 1.3 y 1.7 se descomponen en 41 paquetes técnicos: 18 de plataforma, 10 de seguridad y 13 de servicios comunes. Su estimación ascendente suma 2.024 HH: 776, 512 y 736, respectivamente. El T-14 presenta las fichas con contenido, exclusiones, aceptación, responsable último, esfuerzo por perfil, costo, recursos, supuestos, dependencias e hitos. El T-15 registra las actividades derivadas y compara su carga con las cuentas ya cuantificadas.

La plataforma incluye DEV, QA, PREPROD, PROD y DR: los cinco deberán estar habilitados y operativos para H3, conforme a RT-04.01. Cada ambiente recibe una ficha y evidencia propia; el expediente de entrega consolida esas evidencias sin volver a estimar su ejecución. La creación reproducible y el gobierno del repositorio del cliente se mantienen separados de la construcción de módulos, mientras las pruebas integrales de continuidad y rendimiento pertenecen a 1.10 \parencite[RT-03.03, RT-04.01]{bases-transversales}.

La seguridad comprende identidad central y autoridad local, permisos, claves, consentimiento, auditoría, amenazas, datos de prueba y respuesta técnica. Servicios comunes comprende contratos/eventos, mensajería, sincronización, telemetría, canales y adaptadores externos. Los nueve dominios funcionales consumirán esos componentes sin duplicar su construcción. Las reglas, pantallas y decisiones de cada dominio se estimarán en 1.4--1.6, y las extensiones incrementales conservarán sus paquetes de 1.9.

\subsection{Habilitaciones externas y lotes de terminales}
Los identificadores EXT representan hitos de dependencia, sin esfuerzo de Synaptix adicional al ya incluido en las fichas: EXT-RECINTO, obra y condiciones de instalación recibidas; EXT-EQUIPOS, activos nuevos y conformes a T-11; EXT-RED, tendidos, accesos y enlaces habilitados; EXT-TELEMETRIA, buses y concentradores instalados y accesibles; EXT-LICENCIAS, licencias efectivas y soporte autorizados; EXT-IDENTIDAD, fuente laboral y responsables del cliente; EXT-CONTABLE y EXT-PAGOS, acceso y contratos de integración de terceros. El responsable de coordinación será Fernando Guerrero, con Proyecto y la contraparte designada por el cliente para compras/obras. La coordinación de terceros no convierte sus trabajos en instalaciones ya ejecutadas ni autoriza modificar el alcance o el plazo.

Las habilitaciones de recinto, activos, red y telemetría deben permitir completar las fichas técnicas antes de H3; la fuente de identidad debe estar disponible antes de configurar su integración. Las interfaces contable y de pagos se requieren antes de construir y validar sus adaptadores E2. Proyecto registrará fecha requerida, fecha ofrecida, responsable externo, condición de recepción y efecto sobre dependencias. Si una condición falla, el expediente registrará el bloqueo y se gestionará su resolución dentro de los hitos contractuales.

\paragraph{1.2.3.6.t.b: parque de terminales enrolado y comprobado}
La ficha de referencia 1.2.3.4 define y prueba políticas; este paquete aplica esas políticas al parque. $t$ identifica Windows o Android y $b$ un lote de hasta cuatro equipos, con activo y usuario o función asignados. Bruno Toro responde por el resultado. Incluye enrolamiento, cifrado, versión, actualización, protección, retiro y registro de acceso; excluye adquisición, obras físicas, capacitación de procesos e implementación de la aplicación.

Se recibirá con un registro por activo y comprobación de política, cobertura de protección, bloqueo y recuperación autorizada; un equipo sin licencia o agente efectivo no se considerará protegido. Esfuerzo por lote: Operación 8, Seguridad 4 y Calidad 4 HH, total 16. Recursos: terminales recibidos de T-11, UEM/licencias efectivas, política recibida, red y usuarios autorizados. Se supone hardware compatible y una configuración estándar por tipo; incompatibilidades se dimensionan aparte antes de enrolar.

Para cantidades efectivas $N_W,N_A$, se genera $L=\lceil N_W/4\rceil+\lceil N_A/4\rceil$ lotes, con cero lotes sólo para un tipo confirmado sin equipos. El esfuerzo adicional es $16L$ HH; no se cuenta otra vez la prueba de referencia. El inventario actualizado deberá instanciar los lotes sin reemplazar el detalle de T-11. Cada lote se controla dentro de una quincena y antes de la habilitación del grupo de usuarios que lo necesita; las fechas se fijan con la nivelación y las ventanas autorizadas. La ausencia de inventario confirmado no representa cero horas. Predecesoras: 1.2.3.4, terminales y licencias recibidos; hito: preparación técnica H3 y habilitaciones E1/E2 correspondientes. Referencias: SD4, puestos y gestión de terminales; T-11; SD2, SUP-04; BT capítulo 8.
